\section{Comparación de resultados con una compañera}

En esta sección se registrarán únicamente los requisitos en los que existan diferencias entre la clasificación individual y la realizada por el compañero Alvarez Rai. No es necesario incluir aquellos en los que ambas clasificaciones coincidan.

\begingroup
\small
\setlength{\tw}{\dimexpr\linewidth-8\tabcolsep-5\arrayrulewidth\relax}
\renewcommand{\arraystretch}{2.2}
\begin{longtable}{|L{0.12\tw}|L{0.25\tw}|L{0.25\tw}|L{0.38\tw}|}
\caption{Requisitos con diferencias entre la clasificación propia y la de la compañera.}\label{tab:comparacion}\\
\hline
\filaencabezado \encabezado{Requisito} & \encabezado{Clasificación propia} & \encabezado{Clasificación del compañero Alvarez Rai} & \encabezado{Observación} \\ \hline
\endfirsthead
\hline
\filaencabezado \encabezado{Requisito} & \encabezado{Clasificación propia} & \encabezado{Clasificación de la compañera} & \encabezado{Observación} \\ \hline
\endhead
E10&Calidad: capacidad de interacción – inclusividad y operabilidad. &Calidad: capacidad de interacción – learnability y operabilidad. &Se coincide en la subcaracterística de operabilidad, pero existe una diferencia en la segunda subcaracterística. Se considera la inclusividad por las necesidades de los pacientes adultos mayores, mientras que la otra clasificación contempla la facilidad de aprendizaje (learnability). \\ \hline
E19&Calidad: fiabilidad – recuperabilidad. &Calidad: fiabilidad – recuperabilidad y tolerancia a fallos. &Se coincide en la recuperabilidad, pero existe una diferencia respecto a la inclusión de la tolerancia a fallos. Esta última también puede relacionarse con el enunciado, ya que contempla el fallo del servidor principal y la continuidad del servicio. \\ \hline
\end{longtable}
\endgroup

\subsection{Resumen de la comparación}

Al comparar ambas clasificaciones, se encontraron diferencias únicamente en los requisitos E10 y E19. En el E10, la diferencia corresponde a la subcaracterística de capacidad de interacción considerada; en el E19, se relaciona con la inclusión de la tolerancia a fallos junto con la recuperabilidad. En los otros 18 requisitos existe coincidencia en las clasificaciones realizadas.
